The Cosmic Horseshoe is a gravitationally lensed galaxy at $z=2.379$. It has a magnification factor $\mu=30$. The magnification that occurs due to gravitational lensing is analogous to that provided by an optical telescope, but is due to the gravitational effects of massive bodies. Answer some questions about this galaxy to help a researcher.
\textbf{Note.} AB magnitudes are converted analogously to the magnitude scale that you are familiar with. The average flux density of an object is the flux $F$ of that object divided by the bandwidth $\Delta\nu$ of the detector.
\begin{enumerate}
\item At what wavelength $\lambda_{\mathrm{obs}}$ would you expect to detect the H$\alpha$ emission line emitted by this galaxy? Use $\lambda_{\mathrm{rest}}=656.3\,\mathrm{nm}$.
\item Source plane reconstruction is the technique by which the lensing effects on the galaxy are estimated, in order to get an understanding of the galaxy's morphology. We use the technique and now need to determine the object's physical properties. Using the Hubble Space Telescope (HST), the F606W band AB magnitude is found to be 20.5. What is the object's magnitude after correcting for lensing?
\item One can convert from AB magnitudes to flux density by using the fact that a source at 0 AB magnitude has $3631\,\mathrm{Jy}$ of flux density. What is (a) the apparent flux density of the Cosmic Horseshoe and (b) the magnification corrected flux density?
\item The Cosmic Horseshoe is observed to have its F606W UV continuum emission spread across an area of $13\,\mathrm{arcsec^2}$. What is its apparent UV continuum emission surface density, in units of $\mathrm{mag\,arcsec^{-2}}$?
\item If the angular scale at the Horseshoe's redshift is $8\,\mathrm{kpc\,arcsec^{-1}}$, what is the radius of the galaxy, assuming that it is a face-on disk?
\end{enumerate}
